\chapter{Proposta} \label{cap:Proposta}

% Obs: MM = Monitoramento de Memória

% ------------------------------------

% 1. (Prof. Paulo) Como será feita a coleta das métricas dos servidores?
% R. A coleta será feita por meio de requisições ao sistema operacional, diretamente.

% 2. (Prof. Paulo) Qual será a frequência da coleta?
% R. Foram definidos três períodos distintos de coleta e atualização dos dados - 6s, 60s (1 min) e 600s (10min), com o propósito de avaliar, em diferentes ordens de grandeza, o impacto da desatualização de dados no monitoramento.

% 3. (Prof. Paulo) Como será o monitoramento - Que comando será utilizado? Será criado um programa que coleta o uso de memória?
% R. Nos servidores, haverá uma aplicação que periodicamente usa o sistema operacional para obter dados de consumo de memória (conforme as frequências de teste definidas anteriormente).

% 4. (Prof. Paulo) Como o balanceador de carga obterá as métricas?
% R. A própria aplicação de balanceamento de carga estará inscrita (subscriber) em um \textit{broker} MQTT que periodicamente receberá eventos de atualização dos dados. Reagindo à esse evento, o balanceador de carga atualiza seus dados periodicamente.

% 5. (Prof. Paulo) Como as métricas serão utilizadas no critério de decisão do balanceador de carga (algoritmo)?
% R. A ser definido (memória x disco), interpretação de dados desatualizados;

% 6. (Prof. Paulo) Como serão os cenários de testes?
% R. Análise perceptiva do travamento de vídeos (QoE final) usando \textit{round-robin}, seleção aleatória, MM 6s, MM 60s, MM 600s (todos com técnicas de interpretação de dados desatualizados).

% 7. (Prof. Paulo) Qual será a topologia dos testes?
% R. Apenas uma rede virtual containerizada com alguns servidores e uma aplicação balanceadora de carga. O balanceamento será feito à nível de aplicação.

% 8. (Prof. Paulo) Um balanceador de carga será instrumentado para alcançar os objetivos?
% R. Uma aplicação simples será implementada para esse propósito. 

% 9. (Prof. Paulo) A coleta da informação pelo balanceador de carga e a tomada de decisão irão atrasar em quanto tempo o encaminhamento dos pacotes? Esse atraso irá impactar na qualidade?
% R. A escolha do protocolo MQTT foi feita justamente para este propósito, por ser leve. Espera-se que o impacto maior seja nos período menor (6s), uma vez que demandará maior tempo de processamento do balanceador de carga.  

% 10. (Prof. Paulo) O monitoramento de memória irá aumentar o consumo de outros recursos além da própria memória principal e secundária?
% R. De maneira análoga, espera-se que apenas que o consumo de CPU seja aumentado e com mais intensidade no período menor (6s) visto que as chamadas às APIs do sistema operacional serão mais frequentes.

% Adicional: Colocar aqui, na metodologia (diluído em outras partes) o referencial teórico do artigo "Interpreting Stale Load Information" de modo a obter respostas referenciadas para as perguntas 9 e 10 e fundamentar o algoritmo utilizado (pergunta 5).

% ------------------------------------

Esta seção aborda pormenorizadamente os caminhos escolhidos para a implementação do modelo de balanceamento de carga, bem como os métodos que serão utilizados para realizar a avaliação dos resultados finais.

% (OE 6) Estabelecer, conceitualmente, algoritmo de escolha de servidores por parte do balanceador de carga com base nos dados históricos fornecidos
\section{O algoritmo}

O primeiro passo, crucial para a determinação do modelo de otimização de balanceamento de carga, é a definição do algoritmo de balanceamento. Basicamente, a tarefa principal de um algoritmo de balanceamento de carga é o modo como a escolha do servidor será feita a cada requisição que chega ao sistema balanceador de carga.

Em particular, o algoritmo escolhido será \textbf{probabilístico}. Isso significa que, para cada servidor que poderá ser escolhido no processo, haverá uma probabilidade intrínseca baseada em parâmetros tanto individuais quanto gerais, e a estratégia de determinação dessa probabilidade é o que definirá o algoritmo em si. De acordo com o objetivo principal do trabalho, as duas principais (mas não únicas) métricas que exercerão influência sobre a determinação da probabilidade de escolha de cada servidor são o \textbf{consumo de memória RAM} e a \textbf{Taxa de Leitura de Disco}. Essas duas métricas são essenciais dentro do contexto de servidores que distribuem conteúdo, uma vez que sua única tarefa é recuperar estes arquivos do disco e efetivamente entregá-los ao solicitante. 

Em paralelo, os servidores enviarão, periodicamente, os dados de monitoramento referentes ao consumo de memória principal e à taxa de leitura de disco. A ideia é que o balanceador mantenha uma estrutura dinâmica que será atualizada de forma contínua, atualizando os parâmetros necessários e efetivamente promovendo eficiência na distribuição das requisições. O método que será usado para realizar a implementação deste conceito, bem como a definição das taxas de atualização dos dados dos servidores para o balanceador de carga, serão discutidos posteriormente.

A Equação \ref{eq:basic_li} utilizada no algoritmo corresponde a uma versão adaptada do \textit{Basic LI} \cite{interpretingStaleLoadInformation2000} de modo que, dada uma requisição, a determinação da probabilidade de escolha de um servidor $P{i}$ será determinada de acordo com seus parâmetros, onde:

\begin{enumerate}
    %\item[$T$]: Tempo médio de desatualização da informação de carga, em unidades de tempo de serviço por tarefa.
    \item[$n$]: Número total de servidores.
    %\item[$\lambda$]: Taxa média de chegada de requisições por servidor, como uma fração da vazão máxima do servidor. 
    \item[$L_{i}$]: Carga reportada no servidor $i$ (cálculo determinado na Equação \ref{eq:limemory}). 
    \item[$L_{tot}$]: Soma de todas as cargas reportadas ($\sum L_i$). 
    \item[$arrive_T$]: Informação de carga total esperada para o tempo de vida médio das informações de carga dos servidores (cálculo determinado na Equação \ref{eq:arrive_t}).
    \item[$P_i$]: Probabilidade de que uma requisição que chega seja enviada para o servidor $i$. 
\end{enumerate}

\begin{gather} \label{eq:basic_li}
P_{i}=\frac{\frac{L_{tot}+arrive_{T}}{n}-L_{i}}{arrive_{T}}
\end{gather}

Nota-se que $arrive_T$ é um parâmetro central deste processo. Ele cresce proporcionalmente ao tempo médio de desatualização das informações e, de acordo com a equação anterior (\ref{eq:basic_li}), um valor grande de $arrive_T$ fará com que a distribuição de probabilidade seja uniforme, efetivamente diminuindo a informação de carga. Esse comportamento é esperado e otimizado porque, quanto mais antigas as informações de carga são, menos o algoritmo as considera no balanceamento, que apresentará comportamento similar à seleção aleatória (sem pesos). Todos os parâmetros envolvidos são mantidos e atualizados iterativamente no estado interno do balanceador de carga e são atualizados mediante à chegada de novas informações de carga. O cálculo de $arrive_T$ se dá da seguinte forma:

\begin{enumerate}
\item[$arrive_T$]: Informação de carga total esperada para o tempo de vida médio das informações de carga dos servidores;
\item[$R_{total}$]: Número total de requisições recebidas desde o início do processo da aplicação do balanceador de carga;
\item[$T_{elapsed}$]: Tempo, em segundos, desde o início do processo da aplicação do balanceador de carga;
\item[$T_{stale}$]: Tempo médio de desatualização de dados de carga de todos os servidores;
\item[$L_{tot}$]: Soma de todas as cargas reportadas;
\item[$R_{active}$]: Número total de requisições ativas em todos os servidores, calculado como $R_{active} = \sum_{i=1}^{n} AR_i$, onde $AR_i$ representa o número de requisições ativas no servidor $i$ no momento da coleta da métrica;
\item[$C$]: Coeficiente de suavização para relativização de desatualização alta e baixa (definido heuristicamente mediante experimentação);
\end{enumerate}

\begin{gather} \label{eq:arrive_t} 
arrive_T = \left( \frac{R_{total}}{T_{elapsed}} \right) \cdot T_{stale} \cdot \left( \frac{L_{tot}}{R_{active}} \right) \cdot C \end{gather}

Note-se que no contexto do trabalho \cite{interpretingStaleLoadInformation2000}, "carga" é um termo que remete ao tamanho da fila de requisições em estado de processamento pendente. Para fins de implementação e aderência com o problema proposto por este trabalho, o parâmetro de carga ($L_i$ e consequentemente $L_{tot}$), aplicado na Equação \mbox{\ref{eq:basic_li}}, será calculado previamente através da coleta periódica das métricas de consumo de memória RAM e taxa de leitura de disco (cujo processo será descrito posteriormente), da normalização desses parâmetros em pontos percentuais e, finalmente, da mesclagem desses parâmetros para que possam ser inseridos na fórmula. Nesse sentido, a carga de um servidor ($L_i$) será determinada considerando:

\begin{enumerate}
    \item[$n$]: Número do total de servidores.
    \item[$LM_i$]: Consumo de memória principal no servidor $i$, representando a quantidade de memória usada para transferir o conteúdo no tempo T, normalizado em pontos percentuais com relação à capacidade total de memória do servidor $i$.
    \item[$LD_i$]: Carga do disco no servidor $i$, representando a velocidade do fluxo de busca e leitura das informações do disco no tempo T, normalizada em pontos percentuais com relação a capacidade total de leitura do disco no servidor $i$.
    \item[$TM$]: Taxa de uso geral de memória, determinado por $\displaystyle\sum_{i=0}^n \frac{LM_i}{n}$.
    \item[$TD$]: Taxa de consumo geral de leitura de disco, determinado por $\displaystyle\sum_{i=0}^n \frac{LD_i}{n}$. 
\end{enumerate}

Sendo assim, a Equação \mbox{\ref{eq:limemory}} define a carga, no contexto deste trabalho, de um servidor de conteúdo.

\begin{gather} \label{eq:limemory}
    L_i = CM * LM_i + CD * LD_i
\end{gather}

Nesse contexto, $CM$ e $CD$ são coeficientes que multiplicam as suas respectivas cargas normalizadas. Eles são definidos de acordo com as Equações \ref{eq:cm} e \ref{eq:cd}, respectivamente, e representam a fração normalizada de memória e disco totais, respectivamente, em relação ao total conjugado de memória e disco.

\begin{gather} \label{eq:cm}
    CM = \frac{TM}{TM + TD}
\end{gather}

\begin{gather} \label{eq:cd}
    CD = \frac{TD}{TM + TD}
\end{gather}

Os coeficientes $CM$ e $CD$ são valores entre $0$ e $1$ de modo que $CM + CD = 1$. Eles são determinados dessa forma para que o recurso computacional que mais esteja sendo utilizado de modo geral em todos os servidores efetivamente contribua mais para o valor da carga de cada servidor $L_i$, de forma diretamente proporcional à taxa de consumo de cada um dos parâmetros. Finalmente, para que o sistema efetivamente esteja apto para calcular o valor de cada servidor ($L_i$), basta que ele obtenha os valores de $LM_i$ e $LD_i$ de todos os servidores. Com eles, será possível calcular primeiramente $TM$ e $TD$, e, com esses valores, calcular $CM$ e $CD$ dinamicamente.

O funcionamento do algoritmo ocorre em dois momentos distintos e independentes. O primeiro momento corresponde à \textbf{atualização da distribuição de probabilidade}, que é acionada sempre que uma nova informação de carga chega ao balanceador via MQTT. Neste processo, o balanceador atualiza as métricas do servidor que enviou a informação ($LM_i$, $LD_i$, $T_i$, $AR_i$) e, em seguida, recalcula todos os parâmetros globais necessários: as médias de uso de memória ($TM$) e disco ($TD$), o tempo médio de desatualização ($T$), os coeficientes de ponderação ($CM$ e $CD$), a carga de cada servidor ($L_i$), a carga total ($L_{tot}$), o parâmetro de chegadas esperadas ($arrive_T$) e, finalmente, a probabilidade de seleção de cada servidor ($P_i$). É importante notar que, mesmo quando apenas um servidor envia novas métricas, todos os parâmetros são recalculados considerando todos os servidores conhecidos, garantindo que a distribuição de probabilidade seja sempre consistente e atualizada.

O segundo momento corresponde à \textbf{seleção do servidor}, que ocorre a cada requisição que chega ao balanceador. Neste processo, o balanceador simplesmente utiliza a distribuição de probabilidade mais recente (calculada no primeiro momento) para realizar uma seleção aleatória com pesos. Especificamente, o algoritmo incrementa o contador total de requisições ($R$), obtém um snapshot das probabilidades $P_i$ já calculadas e, então, realiza uma seleção aleatória ponderada, onde cada servidor tem uma chance de ser escolhido proporcional à sua probabilidade $P_i$. Caso nenhuma distribuição de probabilidade tenha sido calculada ainda (situação inicial antes da chegada da primeira métrica), o algoritmo utiliza uma distribuição uniforme, onde todos os servidores têm igual probabilidade de serem escolhidos.

\begin{algorithm}[H]
\caption{Atualização da distribuição de probabilidade}
\begin{algorithmic}[1]
\Statex \textbf{Gatilho}: Chegada de uma nova informação de carga via MQTT.
\Require
    \Statex Para cada servidor $S_i$, suas métricas: $LM_i$, $LD_i$, $T_i$, $AR_i$.
\Ensure
    \Statex Distribuição de probabilidade $P_i$ atualizada para todos os servidores.

\Procedure{AtualizarDistribuicao}{$LM_i, LD_i, T_i, AR_i$}
    \State Atualizar métricas do servidor $S_i$
    
    \State Calcular agregados globais:
    \State $TM \leftarrow \frac{1}{n} \sum_{i=1}^{n} LM_i$, $TD \leftarrow \frac{1}{n} \sum_{i=1}^{n} LD_i$
    \State $T \leftarrow \frac{1}{n} \sum_{i=1}^{n} (T_{atual} - T_i)$ \Comment{Tempo médio de desatualização}
    
    \State Calcular coeficientes $CM$ e $CD$ conforme Equações \ref{eq:cm} e \ref{eq:cd}
    
    \State Calcular carga $L_i$ para cada servidor conforme Equação \ref{eq:limemory}
    \State $L_{tot} \leftarrow \sum_{i=1}^{n} L_i$
    
    \State Calcular $arrive_T$ conforme Equação \ref{eq:arrive_t}
    
    \State Calcular probabilidade $P_i$ para cada servidor conforme Equação \ref{eq:basic_li}
    \If{nenhum servidor disponível \textbf{ou} $arrive_T \leq 0$}
        \State $P_i \leftarrow 1/n$ para todos os servidores \Comment{Distribuição uniforme}
    \EndIf
\EndProcedure
\end{algorithmic}
\end{algorithm}

Para o funcionamento do algoritmo, o balanceador de carga mantém em memória um conjunto de estruturas de dados que armazenam o estado atual do sistema. Essas estruturas incluem: as métricas individuais de cada servidor ($LM_i$, $LD_i$, $T_i$, $AR_i$), as cargas calculadas ($L_i$), as probabilidades de seleção ($P_i$), além de variáveis agregadas como $L_{tot}$ e $n$. Todas essas informações são dinamicamente atualizadas sempre que uma nova métrica chega via MQTT, e o número de servidores ($n$) é automaticamente ajustado quando novos servidores são detectados através da identificação única presente nas mensagens MQTT. A detecção de novos servidores ocorre de forma automática: quando uma métrica é recebida de um servidor ainda não conhecido, ele é automaticamente adicionado ao conjunto de servidores monitorados.

% (OE 7) Definição conceitual da arquitetura e definição dos requisitos topológicos da rede e da estrutura dos servidores que definirão o contexto da aplicação da metodologia
% (OE 5) Definição da estratégia escolhida de transmissão contínua de dados de monitoramento de servidores de distribuição de conteúdo para o serviço balanceador de carga
\section{Os requisitos tecnológicos}

A infraestrutura experimental para este trabalho foi implementada em um ambiente virtualizado, contido integralmente em uma única máquina física. A escolha por esta abordagem, suportada pela tecnologia de \textit{contêineres} \textit{Docker}, é estratégica e visa mitigar a influência de variáveis de rede externas, como latência e instabilidade de conexão. Ao confinar toda a comunicação a uma rede virtual interna gerenciada pelo \textit{Docker}, o foco da análise é direcionado exclusivamente ao impacto do monitoramento de recursos, como a memória, no desempenho e na otimização do balanceador de carga, que é o objetivo central deste trabalho. Adicionalmente, a utilização do \textit{Docker} como plataforma padrão para todos os componentes do sistema garante a padronização do ambiente de execução, eliminando problemas de dependência de sistemas operacionais e bibliotecas, o que assegura a consistência e a reprodutibilidade dos experimentos.

Neste ecossistema controlado, quatro aplicações principais interagem para formar o sistema em estudo: As três principais são o servidor de conteúdo, que será replicado de modo redundante para que o cenário emule simplificadamente uma CDN, o cliente e o balanceador de carga. Em conjunto, ainda haverá mais um serviço auxiliar que fará interações diretamente com os servidores e com o balanceador de carga, sendo ele o \textit{broker} de mensagens MQTT. O servidor de conteúdo será responsável por fornecer segmentos de vídeo de acordo com o padrão MPEG-DASH. O balanceador de carga atuará como um intermediário, recebendo as requisições dos clientes e distribuindo-as de forma inteligente, baseado na avaliação de carga, entre os servidores disponíveis. A aplicação cliente iniciará várias requisições de modo paralelo, direcionando-as ao balanceador de carga (solicitando segmentos subsequentes conforme o MPEG-DASH) e, subsequentemente, consumindo o conteúdo de mídia retornado, fechando assim o ciclo de interação do sistema. Em paralelo, os próprios servidores, periodicamente, fornecerão seus dados de carga de memória principal e secundária publicando-os em tópicos específicos do \textit{broker} por meio do modelo de comunicação indireta \textit{publish-subscribe} \cite{eugster2003many}. O balanceador de carga, que estará inscrito como assinante desses tópicos, eventualmente consumirá esses dados de modo a realimentar os parâmetros internos utilizados para os cálculos envolvidos (ver Equações \ref{eq:basic_li} e \ref{eq:limemory}).

A arquitetura do sistema foi dimensionada para simular um cenário de alta demanda e é composta por uma variedade de serviços. O conjunto de recursos consiste em 8 \textit{containers} atuando como servidores MPEG-DASH, constituindo o \textit{pool} de nós para os quais a carga será distribuída. Há um contêiner dedicado para atuar como \textit{broker} MQTT, com o propósito de mediar a comunicação publicar-assinar para a atualização dos parâmetros do balanceador de carga. Também haverá um serviço de monitoramento com acesso direto tanto aos servidores quanto ao \textit{broker} MQTT sob a finalidade de auxiliar no processo de desenvolvimento e depuração, fornecendo observabilidade. O sistema contará com um único balanceador de carga, que representa a peça central da investigação. Para gerar uma carga de trabalho realista e concorrente, será implementado um \textit{script} cliente simples, projetado para gerar múltiplas requisições paralelas, simulando o comportamento de um grande número de usuários acessando o serviço simultaneamente por meio de um \textit{player} MPEG-DASH. Finalmente, um simples \textit{player} de vídeo baseado em navegador será utilizado como ferramenta de validação, cujo propósito é permitir uma avaliação qualitativa da experiência do usuário (QoE) ao reproduzir o vídeo disposto pelos servidores, mesmo sob a sobrecarga imposta pela aplicação cliente de teste.

A seleção de tecnologias para cada componente foi alinhada aos objetivos do projeto. Duas imagens \textit{Docker} pré-existentes serão configuradas e usadas para criar os servidores MPEG-DASH e o \textit{broker} MQTT (\textit{Message Queuing Telemetry Transport}), respectivamente. Uma aplicação será desenvolvida e encapsulada em um contêiner na forma de um \textit{script} \textit{Python}, de modo que será responsável por simular sobrecarga nos serviços e no balanceador de carga. Por fim, o balanceador de carga, componente central desta pesquisa, será desenvolvido com \textit{design} modular e aberto, permitindo a fácil substituição e modificação da estratégia de balanceamento, o que viabilizará futuras comparações de desempenho com algoritmos generalistas, como \textit{Round-Robin} e Seleção Aleatória.

De modo abstrato, a Figura \ref{fig:arch} apresenta o panorama geral da arquitetura final do sistema criado de acordo com a Figura \mbox{\ref{fig:arch}}.

\begin{figure}[htbp]
    \caption{Arquitetura geral do sistema.}
    \centering
    \includegraphics[width=1\linewidth]{./Textuais/3-Proposta/arch.png}
    \text{Fonte: Elaborado pelo autor (2025)}
    \label{fig:arch}
\end{figure}

\section{Considerações}

Este capítulo apresentou a proposta metodológica para a implementação de um sistema de balanceamento de carga baseado em monitoramento de recursos computacionais, especificamente memória RAM e taxa de leitura de disco. O algoritmo proposto, fundamentado no trabalho de \cite{interpretingStaleLoadInformation2000}, utiliza uma abordagem probabilística que adapta dinamicamente a distribuição de requisições entre servidores conforme o estado atual de carga de cada nó.

A escolha por métricas de memória e disco justifica-se pelo contexto de aplicação: servidores de distribuição de conteúdo que recuperam arquivos do disco e os transferem para clientes, sendo esses recursos críticos para o desempenho do sistema. A implementação do algoritmo em dois momentos distintos --- atualização da distribuição de probabilidade e seleção de servidor --- permite que o balanceador mantenha informações atualizadas sobre o estado do sistema sem impactar significativamente a latência de processamento de requisições.

A arquitetura proposta, baseada em contêineres \textit{Docker} e comunicação via MQTT, oferece um ambiente controlado e reproduzível para a avaliação do algoritmo, isolando variáveis externas e permitindo que o foco da análise seja direcionado exclusivamente ao impacto do monitoramento de recursos no balanceamento de carga. A modularidade do sistema, com estratégias de balanceamento intercambiáveis, viabiliza comparações futuras com algoritmos generalistas e permite a extensão da pesquisa para outros contextos de aplicação.

Os próximos capítulos detalharão a implementação prática desta proposta, os experimentos realizados e os resultados obtidos, fornecendo evidências empíricas sobre a eficácia da abordagem proposta e seu impacto na qualidade de experiência dos usuários finais.